\chapter{Tiny210
设备模型与驱动开发}\label{tiny210-ux8bbeux5907ux6a21ux578bux4e0eux9a71ux52a8ux5f00ux53d1}

\section{分层}\label{ux5206ux5c42}

当前代码借鉴 Linux 的 bus/device/driver 和 initcall 思想，但保持 xv6
的小型实现：

\begin{figure}[htbp]
\centering
\begin{tikzpicture}[node distance=7mm]
  \node[hbox] (board) {\texttt{devices.c}\\platform device + resource};
  \node[box,right=18mm of board] (drv) {initcall\\\texttt{platform\_driver\_register()}};
  \coordinate (mid) at ($(board)!0.5!(drv)$);
  \node[box,below=of mid] (bus) {platform bus\\按 name 匹配};
  \node[box,below=of bus] (probe) {\texttt{driver->probe(device)}\\MMIO、clock、IRQ、DMA};
  \node[hbox,below left=7mm and 10mm of probe] (chr) {字符设备\\\texttt{file\_operations}};
  \node[hbox,below right=7mm and 10mm of probe] (net) {网络设备\\\texttt{net\_device}};
  \draw[arr] (board)--(bus); \draw[arr] (drv)--(bus); \draw[arr] (bus)--(probe);
  \draw[arr] (probe)--(chr); \draw[arr] (probe)--(net);
\end{tikzpicture}
\caption{Tiny210 的 platform device/driver 分层与绑定}
\label{fig:tiny-driver-model}
\end{figure}


板级代码只描述``有什么硬件和资源''；驱动只描述``如何操作这种控制器''。

\section{核心对象}\label{ux6838ux5fc3ux5bf9ux8c61}

定义位于
\href{../../../kernel/device.h}{\passthrough{\lstinline!kernel/device.h!}}：

\begin{lstlisting}[language=C]
struct resource {
  uint64 start, end;
  int type;                 // IORESOURCE_MEM / IORESOURCE_IRQ
  char *name;
};

struct device {
  char *name;
  int id;
  struct bus_type *bus;
  struct resource resource[DEVICE_RES_MAX];
  int nresource;
  struct device_driver *driver;
  void *driver_data;        // 板级参数或 probe 后私有数据
};

struct device_driver {
  char *name;
  struct bus_type *bus;
  int (*probe)(struct device *dev);
  void (*remove)(struct device *dev);
};
\end{lstlisting}

\passthrough{\lstinline!platform\_bus.match()!} 当前只比较
\passthrough{\lstinline!dev->name == drv->name!}。没有 device tree
compatible、ACPI、modalias 或延迟 probe。

\section{注册和匹配}\label{ux6ce8ux518cux548cux5339ux914d}

启动顺序是：

\begin{lstlisting}
device_init()
  -> 初始化静态 device_core
  -> bus_register(&platform_bus)

s5pv210_register_devices()
  -> platform_device_register(clock/wdt/vic/pwm/uart/gpio/fb)
  -> 此时大多数 driver 尚未注册，device 暂时 unbound

do_initcalls()
  -> 各驱动 platform_driver_register()
  -> driver_register() 遍历已存在 device
  -> name 匹配
  -> bus_probe()
  -> driver->probe(dev)
\end{lstlisting}

反过来，如果 driver 先存在，\passthrough{\lstinline!device\_register()!}
也会遍历驱动并立即匹配。
\href{../../../kernel/device.c}{\passthrough{\lstinline!device.c!}} 支持
device/driver/bus unregister，ARM32 IRQ 层也有
\passthrough{\lstinline!free\_irq()!}。不过多数 S5PV210
外设是不可热拔插的板载设备，驱动的 \passthrough{\lstinline!.remove!}
和资源 回滚仍不完整，不能把当前模型视为完整热拔插框架。

\section{initcall}\label{initcall}

\href{../../../kernel/module.h}{\passthrough{\lstinline!module.h!}}
把静态 \passthrough{\lstinline!module\_desc!} 放入链接器
\passthrough{\lstinline!.initcall.list!}，\href{../../../kernel/arch/arm/initcall.c}{\passthrough{\lstinline!initcall.c!}}
按 \passthrough{\lstinline!(level, order)!} 稳定排序执行：

{\def\LTcaptype{none} % do not increment counter
\begin{longtable}[]{@{}p{0.18\linewidth}p{0.26\linewidth}p{0.50\linewidth}@{}}
\toprule\noalign{}
level & 宏 & Tiny210 用途 \\
\midrule\noalign{}
\endhead
\bottomrule\noalign{}
\endlastfoot
1 core & \passthrough{\lstinline!core\_initcall!} & clock、关闭
watchdog \\
3 arch & \passthrough{\lstinline!arch\_initcall!} & 根 VIC \\
4 subsys & \passthrough{\lstinline!subsys\_initcall!} & PWM
clocksource/tick、input core \\
5 fs & \passthrough{\lstinline!fs\_initcall!} & 当前 Tiny210
无内置条目 \\
6 device & \passthrough{\lstinline!module\_init!} &
serial、K1-\/-K8、LCD；\passthrough{\lstinline!s3c-fb!} 用 order 1 \\
7 late & \passthrough{\lstinline!late\_initcall!} & 当前 Tiny210
无内置条目 \\
\end{longtable}
}

典型依赖关系是 clock → IRQ controller → timer → 普通设备。相同 level
有明确依赖时 使用
\passthrough{\lstinline!module\_init\_ord()!}，不要依赖目标文件链接顺序。

\section{当前 S5PV210 platform
devices}\label{ux5f53ux524d-s5pv210-platform-devices}

{\def\LTcaptype{none} % do not increment counter
\begin{longtable}[]{@{}p{0.15\linewidth}p{0.20\linewidth}p{0.20\linewidth}p{0.35\linewidth}@{}}
\toprule\noalign{}
device/driver name & resource & initcall & 当前行为 \\
\midrule\noalign{}
\endhead
\bottomrule\noalign{}
\endlastfoot
\passthrough{\lstinline!s5pv210-clock!} & clock MMIO
\passthrough{\lstinline!0xE0100000!} & core & 读取 APLL/MPLL
和总线频率 \\
\passthrough{\lstinline!s3c2410-wdt!} & WDT
\passthrough{\lstinline!0xE2700000!} & core &
\passthrough{\lstinline!WTCON=0!}，只关闭看门狗 \\
\passthrough{\lstinline!s5pv210-vic!} & VIC0..3
\passthrough{\lstinline!0xF2000000!} & arch & 安装 IRQ 0..127
根分派器 \\
\passthrough{\lstinline!s5pv210-pwm!} & PWM
\passthrough{\lstinline!0xE2500000!}, IRQ25 & subsys & Timer 4
tick/clocksource \\
\passthrough{\lstinline!s5pv210-uart!} & UART0, IRQ42 & device &
控制台；RX 当前按 tick 轮询 \\
\passthrough{\lstinline!s5pv210-gpio!} & GPIO
\passthrough{\lstinline!0xE0200000!} & --- & 只有 device，当前无匹配
driver \\
\passthrough{\lstinline!s5pv210-keys!} & GPIO/EINT, IRQ16 & device &
K1-\/-K8，中断 + 20 ms 延迟消抖，输出到
\passthrough{\lstinline!event0!} \\
\passthrough{\lstinline!s3c-fb!} & FIMD + framebuffer & device order 1 &
Window0 XRGB8888 + fbcon \\
\end{longtable}
}

此外 \passthrough{\lstinline!ramdisk\_embed.c!}
是构建时直接链接的块后端，没有 platform device/driver 和 initcall；它由
\passthrough{\lstinline!bread()!} 路径直接调用。

\section{MMIO 映射}\label{mmio-ux6620ux5c04}

\passthrough{\lstinline!devices.c!} 的
\passthrough{\lstinline!x210\_map\_devices()!} 在
\passthrough{\lstinline!kvminit()!} 阶段遍历每个
\passthrough{\lstinline!IORESOURCE\_MEM!}：

\begin{itemize}
\tightlist
\item
  普通控制器寄存器：1 MiB section、privileged RW、XN、strongly ordered；
\item
  名称为 \passthrough{\lstinline!framebuffer!} 的资源：privileged
  RW、XN、Normal non-cacheable。
\end{itemize}

因此 probe 中可直接把 \passthrough{\lstinline!resource.start!}
转为指针。当前是恒等映射，尚未实现 Linux 风格的
\passthrough{\lstinline!ioremap()!} 动态虚拟地址；新增资源若不进入 board
device 表，也不会自动 出现在页表中。

\section{字符设备与
file\_operations}\label{ux5b57ux7b26ux8bbeux5907ux4e0e-file_operations}

硬件 driver 如果需要用户接口，可调用：

\begin{lstlisting}[language=C]
register_chrdev(major, "name", &fops);
\end{lstlisting}

\href{../../../kernel/device.h}{\passthrough{\lstinline!file\_operations!}}
支持：

\begin{itemize}
\tightlist
\item
  简单的 \passthrough{\lstinline!read/write!}；
\item
  每次 open 的 \passthrough{\lstinline!open/release!}；
\item
  可看到 \passthrough{\lstinline!struct file!} 和
  \passthrough{\lstinline!private\_data!} 的
  \passthrough{\lstinline!fread!}。
\end{itemize}

文件系统中的 \passthrough{\lstinline!T\_DEVICE!} inode 保存
major/minor。\passthrough{\lstinline!open()!} 后由字符设备表找到 fops，
read/write 再分派到驱动。device core 不会自动创建
\passthrough{\lstinline!/dev!} 节点；Tiny210 的
\passthrough{\lstinline!/init!} 当前使用
\passthrough{\lstinline!mknod()!} 创建 console、ttyS0、fb0 和
\passthrough{\lstinline!input/event0!}。

\section{板载 K1-\/-K8：输入子系统和 event
接口}\label{ux677fux8f7d-k1--k8ux8f93ux5165ux5b50ux7cfbux7edfux548c-event-ux63a5ux53e3}

Tiny210 SDK 底板的按键是低电平有效的 GPIO 外部中断输入：

{\def\LTcaptype{none} % do not increment counter
\begin{longtable}[]{@{}p{0.15\linewidth}p{0.20\linewidth}p{0.20\linewidth}p{0.35\linewidth}@{}}
\toprule\noalign{}
按键 & GPIO & 外部中断 & event code \\
\midrule\noalign{}
\endhead
\bottomrule\noalign{}
\endlastfoot
K1-\/-K4 & GPH2{[}0:3{]} & EINT16-\/-19 &
\passthrough{\lstinline!KEY\_1!}-\/-\passthrough{\lstinline!KEY\_4!} \\
K5-\/-K8 & GPH3{[}0:3{]} & EINT24-\/-27 &
\passthrough{\lstinline!KEY\_5!}-\/-\passthrough{\lstinline!KEY\_8!} \\
\end{longtable}
}

EINT16-\/-31 在 VIC0 上共用 IRQ 16。驱动
\href{../../../kernel/drivers/input/keyboard/s5pv210-keys.c}{\passthrough{\lstinline!s5pv210-keys.c!}}
把 GPH2/GPH3
的前四个引脚配置为外部中断功能、内部上拉、双边沿触发。硬中断只完成三件事：读取
\passthrough{\lstinline!EXT\_INT\_*\_PEND!}、屏蔽并清除触发的按键、把 20
ms 延迟工作加入 \passthrough{\lstinline!system\_wq!}。它不在 IRQ
上下文中忙等。

\begin{lstlisting}
按下或释放 K1--K8（低电平为按下）
  -> GPH2/GPH3 检测双边沿，置 EXT_INT_2/3_PEND
  -> VIC0 IRQ 16
  -> keys_irq(): mask + acknowledge + schedule_delayed_work(20 ms)
  -> keys_debounce_work(): 读取稳定的 GPH2DAT/GPH3DAT，再解除 mask
  -> input_report_key(KEY_n, 1/0)
  -> input_sync(): EV_SYN / SYN_REPORT
  -> evdev_client 环形缓冲
  -> read(/dev/input/event0)
\end{lstlisting}

按一次 K1 会产生一个按下帧和一个释放帧，每帧以
\passthrough{\lstinline!SYN\_REPORT!} 结束：

\begin{lstlisting}
EV_KEY  code=KEY_1(2) value=1
EV_SYN  code=SYN_REPORT value=0
EV_KEY  code=KEY_1(2) value=0
EV_SYN  code=SYN_REPORT value=0
\end{lstlisting}

\passthrough{\lstinline!input.c!} 是设备无关的输入核心。每次打开
\passthrough{\lstinline!event0!} 都分配独立的
\passthrough{\lstinline!evdev\_client!}
环形缓冲，因此多个读者不会互相抢事件；读者落后时会收到
\passthrough{\lstinline!SYN\_DROPPED!}。驱动只调用
\passthrough{\lstinline!input\_report\_key()!}/\passthrough{\lstinline!input\_sync()!}，不直接了解字符设备或用户地址。

实板测试：

\begin{lstlisting}[language=sh]
/bin/evtest
# 或 /bin/evtest /dev/input/event0
\end{lstlisting}

预期启动日志包含：

\begin{lstlisting}
input: Tiny210 GPIO Keys registered as event0
input: Tiny210 K1-K8 GPIO keys on EINT16-19/24-27, IRQ 16
\end{lstlisting}

该接口是事件设备，不会自动把 K1-\/-K8 转成 shell
字符；需要控制菜单或窗口系统的用户 程序读取
\passthrough{\lstinline!event0!} 并解释
\passthrough{\lstinline!KEY\_1!}-\/-\passthrough{\lstinline!KEY\_8!}。

\section{新增一个 platform
driver}\label{ux65b0ux589eux4e00ux4e2a-platform-driver}

\subsection{描述板级资源}\label{ux63cfux8ff0ux677fux7ea7ux8d44ux6e90}

在
\href{../../../kernel/arch/arm/mach-s5pv210/devices.c}{\passthrough{\lstinline!devices.c!}}
增加：

\begin{lstlisting}[language=C]
{ .name = "s5pv210-demo",
  .resource = {
    MEM(0xE1230000, 0x1000, "demo"),
    IRQ(60, "demo")
  },
  .nresource = 2 },
\end{lstlisting}

这一步同时让 \passthrough{\lstinline!x210\_map\_devices()!} 建立 MMIO
映射。

\subsection{编写驱动}\label{ux7f16ux5199ux9a71ux52a8}

\begin{lstlisting}[language=C]
static void demo_irq(int irq, void *arg)
{
  struct device *dev = arg;
  /* 读取并清除设备 pending；不要睡眠 */
}

static int demo_probe(struct device *dev)
{
  uintptr regs = dev->resource[0].start;
  int irq = dev->resource[1].start;
  /* 配置时钟、pinmux 和控制器 */
  return request_irq(irq, demo_irq, dev);
}

static struct device_driver demo_driver = {
  .name = "s5pv210-demo",
  .probe = demo_probe,
};

static void demo_init(void)
{
  platform_driver_register(&demo_driver);
}

module_init("demo", demo_init);
\end{lstlisting}

\subsection{加入构建}\label{ux52a0ux5165ux6784ux5efa}

把不带 \passthrough{\lstinline!.c!} 的相对路径加入
\passthrough{\lstinline!ports/tiny210/Makefile!} 的
\passthrough{\lstinline!DRIVERS!}。若 QEMU 没有对应 设备，只加入非
\passthrough{\lstinline!QEMU\_TEST!} 分支，或者提供一个 QEMU 替代
driver。

\subsection{验证}\label{ux9a8cux8bc1}

\begin{enumerate}
\def\labelenumi{\arabic{enumi}.}
\tightlist
\item
  \passthrough{\lstinline!make tiny210-clean \&\& make tiny210!}，确保交叉编译和链接通过。
\item
  \passthrough{\lstinline!kernel.map!} 中确认 driver/initcall 未被
  \passthrough{\lstinline!--gc-sections!} 删除。
\item
  启动日志确认 \passthrough{\lstinline!initcall: demo!}。
\item
  确认 device 与 driver 名称完全相同且 probe 被调用。
\item
  读回关键寄存器；如果使用 IRQ，确认 pending 清除和中断计数。
\end{enumerate}

\section{当前模型与 Linux
的差异}\label{ux5f53ux524dux6a21ux578bux4e0e-linux-ux7684ux5deeux5f02}

{\def\LTcaptype{none} % do not increment counter
\begin{longtable}[]{@{}p{0.34\linewidth}p{0.60\linewidth}@{}}
\toprule\noalign{}
当前 xv6 & Linux \\
\midrule\noalign{}
\endhead
\bottomrule\noalign{}
\endlastfoot
静态数组，\passthrough{\lstinline!DEVICE\_MAX=16!} &
动态对象/kobject/klist \\
platform 只按 name 匹配 & OF/ACPI/platform id table \\
resource 直接恒等映射 & request\_mem\_region +
ioremap/devm\_ioremap\_resource \\
简单 \passthrough{\lstinline!request\_irq!} &
irqdomain、irqchip、shared/threaded IRQ、devm \\
内置 \passthrough{\lstinline!.initcall.list!} & 内置 initcall + 可装卸
\passthrough{\lstinline!.ko!} module \\
\passthrough{\lstinline!/init!} 手工 mknod & devtmpfs +
uevent/udev/mdev \\
基本 remove 框架 & 完整引用计数、runtime PM、热拔插生命周期 \\
\end{longtable}
}

Tiny210 bring-up
阶段应优先保持路径简单；新增复杂子系统时，应先补生命周期、资源
所有权和错误回滚，而不是只复制 Linux 结构体名称。
